\chapter{Goblinoids — Threshold's Interest}
\label{chap:goblinoids}

\epigraph{"They do not remember where they came from. They only know that they are here now, and that the world asks something of them."}{---Lerris Fenwood, from a field journal found in the Archivolt}

The goblinoids are the Threshold's keepers. They are not fae in the traditional sense—they lack the Trow's ancient courtesy, the Neighbor's patient attention, the Green Court's intricate law. They are \textbf{newer} than the other threshold folk, and they are \textbf{eager}. They remember nothing of their origins. They only know that they exist, that the world is cold, and that survival requires taking what is not freely given.

Goblins, hobgoblins, and bugbears are the three faces of this eagerness. They appear wherever the Threshold's attention pools: abandoned crossroads, the roots of dying trees, the spaces between counted things. They are not evil—they are \textbf{desperate}. But desperation, in a world of balance, is its own kind of strain.

This chapter catalogs the goblinoids as I have encountered them: their habits, their bargains, and the quiet terror of their persistence. They are not as old as the fae, but they are \textbf{more patient}. And patience, as the Threshold teaches, is the most dangerous currency of all.

\section{Goblins — Threshold's Fingers}

\textbf{Description}: Small, sharp-toothed, and quick. Goblins are the most common of the goblinoids, appearing in small bands wherever thresholds are thin. They wear scavenged clothing and carry mismatched tools—a stolen knife, a bent fork, a length of rusted chain. Their eyes are too bright, their grins too wide. They do not blink often.

Goblins are \textbf{fodder}. They are not meant to be a significant threat individually, but in numbers they become a pressure that wears down even the most prepared party. A single goblin is a nuisance. A dozen goblins are a problem. A hundred goblins are a siege.

\subsection{Goblin NPC Template}

\textbf{Tier}: I—II | \textbf{Cap}: 1 | \textbf{Corporeal}, Quick, Opportunistic

\textbf{Dice Pools}:
\begin{itemize}
    \item \textbf{Stealth}: 4d (Wits 2 + Stealth 2)
    \item \textbf{Melee}: 3d (Body 2 + Melee 1) — \textit{Harm 2 (knife or sharpened tool)}
    \item \textbf{Subterfuge}: 4d (Wits 2 + Subterfuge 2)
    \item \textbf{Notice}: 3d (Wits 2 + Notice 1)
\end{itemize}

\textbf{Armor}: None (scavenged cloth, no conversion)

\textbf{Traits}:
\begin{itemize}
    \item \textbf{Threshold's Cunning} — Once per scene, the goblin may reroll a failed Stealth or Subterfuge roll. The GM gains 1 SB.
    \item \textbf{Nimble Escape} — As a Reaction when the goblin would take Harm, it may Disengage and move one range band without provoking Opportunity Attacks.
    \item \textbf{Scavenger's Eye} — Goblins can spot hidden or overlooked objects with a Wits + Notice test (DV 2). Success reveals a useful item, a hidden door, or a clue in the scene.
\end{itemize}

\textbf{Etiquette}: Goblins respect cleverness and fear. A display of quick thinking (a solved puzzle, a spotted trap) grants \textbf{+1 Position} in negotiation. A display of overwhelming force grants \textbf{+1 Position} in intimidation. Goblins do not bargain in good faith—they bargain in \textbf{survival}. Every deal is a test of who is smarter.

\textbf{Complications}: Goblins are loyal to their band, not to any individual. A goblin will betray a rival for a better offer, but it will also avenge a fallen packmate if the blood-price is not paid. They do not forget insults. They do not forgive easily.

\textbf{Resolution}: Most goblins can be bribed, threatened, or tricked into cooperation:
\begin{itemize}
    \item \textbf{Bribe}: Spend 1 Boon to offer food, shiny objects, or a promise of safety. The goblin leaves or provides a minor service.
    \item \textbf{Threaten}: Succeed on a Presence + Command test (DV 3). On a Clean Success, the goblin flees. On a Partial, it flees but takes something on the way out.
    \item \textbf{Trick}: Succeed on a Wits + Subterfuge test (DV 3). On success, the goblin is outsmarted and leaves embarrassed. On a Miss, the GM gains 2 SB and the goblin becomes hostile.
\end{itemize}

\subsection{Goblin Variants}

Goblins are endlessly adaptable. The following variants can be used to add flavor to encounters.

\subsubsection{Tunnel-Goblin (Aeler Underways)}

\textbf{Description}: Pale, blind, and skittering. Tunnel-goblins have adapted to the darkness of the deep ways. They navigate by sound and vibration, and they are \textbf{fast}.

\textbf{Stats}: Tier I, Cap 1. Corporeal, Blind, Fast.\\
\textbf{Harm}: 2. \textbf{Fatigue}: Body 1.

\textbf{Traits}:
\begin{itemize}
    \item \textbf{Echolocation} — Immune to visual stealth. Can detect any creature moving within Near range automatically.
    \item \textbf{Skittering Charge} — Once per scene, the goblin may move two range bands and attack in the same turn without Fatigue.
    \item \textbf{Light Blindness} — In bright light, the goblin suffers \textbf{-1 die} to all actions.
\end{itemize}

\subsubsection{Ember-Goblin (Zahk'Kozahn Fringes)}

\textbf{Description}: Blackened skin, smoldering eyes. Ember-goblins have been touched by the fires that seep from the cursed city's depths. They carry the heat of Malachai's forge.

\textbf{Stats}: Tier I, Cap 1. Corporeal, Smoldering, Unstable.\\
\textbf{Harm}: 2. \textbf{Fatigue}: Body 1.

\textbf{Traits}:
\begin{itemize}
    \item \textbf{Smoldering Touch} — On a successful melee attack, the target marks 1 Fatigue (heat).
    \item \textbf{Ember Burst} — When the goblin is killed, it explodes in a burst of ash and heat. All creatures in Close range must test Body + Athletics (DV 3) or mark 1 Fatigue.
    \item \textbf{Water Vulnerability} — A splash of water deals Harm 2 to the goblin and forces a Resolve test (DV 3) or it flees.
\end{itemize}

\subsubsection{Sea-Goblin (Linn Coast)}

\textbf{Description}: Webbed fingers, seaweed hair, salt-crusted skin. Sea-goblins patrol the tide-lines, scavenging what the sea brings in.

\textbf{Stats}: Tier I, Cap 1. Corporeal, Amphibious, Opportunistic.\\
\textbf{Harm}: 2. \textbf{Fatigue}: Body 1.

\textbf{Traits}:
\begin{itemize}
    \item \textbf{Amphibious} — Can breathe water and move at normal speed while swimming.
    \item \textbf{Tide-Cunning} — While in or near water, the goblin gains \textbf{+1 die} to Stealth and Athletics.
    \item \textbf{Drowning Grasp} — Once per scene, the goblin may attempt to drag a target underwater (Body + Athletics vs DV 3). On success, the target marks 1 Fatigue and is pulled one range band toward the water.
\end{itemize}

\subsubsection{Star-Goblin (High Passes)}

\textbf{Description}: Dust-gray, with eyes that reflect constellations. Star-goblins live in the high passes, where the air is thin and the stars are close.

\textbf{Stats}: Tier I, Cap 1. Corporeal, Star-Touched, Unstable.\\
\textbf{Harm}: 2. \textbf{Fatigue}: Body 1.

\textbf{Traits}:
\begin{itemize}
    \item \textbf{Star-Struck} — Once per scene, the goblin may reroll a failed Notice or Lore roll. The GM gains 1 SB.
    \item \textbf{Starlight Shimmer} — In starlight, the goblin has \textbf{Dominant Position} on Stealth rolls.
    \item \textbf{Star-Fall} — When the goblin is killed, it leaves behind a small shard of crystal that glows faintly in the dark. The shard can be used as a minor light source or sold for 1 Supply.
\end{itemize}

\section{Hobgoblins — Threshold's Fist}

\textbf{Description}: Taller, broader, and more disciplined than their smaller kin. Hobgoblins move with the precision of soldiers, their armor patched and scavenged but well-maintained. They carry weapons with purpose—not the goblin's improvised sharpness, but the blade of a creature that has learned to kill efficiently.

Hobgoblins are the Threshold's \textbf{enforcers}. They appear where order has broken down—where a village has refused to pay its dues, where a road has been closed for too long, where a promise has been broken and not mended. They do not destroy. They \textbf{restore}. By force.

\subsection{Hobgoblin NPC Template}

\textbf{Tier}: II—III | \textbf{Cap}: 3 | \textbf{Corporeal}, Disciplined, Tactical

\textbf{Dice Pools}:
\begin{itemize}
    \item \textbf{Melee}: 6d (Body 3 + Melee 3) — \textit{Harm 4 (sword or polearm)}
    \item \textbf{Command}: 5d (Presence 3 + Command 2)
    \item \textbf{Athletics}: 5d (Body 3 + Athletics 2)
    \item \textbf{Endurance}: 5d (Body 3 + Endurance 2)
\end{itemize}

\textbf{Armor}: Medium (scavenged plate) — \textit{Harm 2 → Fatigue 1, Harm 3 → Fatigue 2}

\textbf{Traits}:
\begin{itemize}
    \item \textbf{Threshold's Discipline} — Once per scene, the hobgoblin may ignore a Fatigue penalty for one roll. The GM gains 1 SB.
    \item \textbf{Tactical Awareness} — Hobgoblins cannot be surprised while their band is intact. They gain \textbf{+1 Position} on the first action of any combat.
    \item \textbf{Oath-Broken Armor} — Hobgoblins are bound by a silent oath to their band. If their band is broken (half or more incapacitated), they suffer \textbf{-1 die} to all actions until they retreat or reform.
\end{itemize}

\textbf{Etiquette}: Hobgoblins respect strength and order. A character who speaks in commands, who holds a line, who does not flinch—these are the ones hobgoblins will deal with. They are not interested in clever tricks. They are interested in \textbf{what works}. A hobgoblin bargain is a transaction: service for service, blood for blood.

\textbf{Complications}: A hobgoblin will never betray its band voluntarily. If a hobgoblin's band is threatened, it will fight to the death. If a hobgoblin's band is destroyed, it will become a lone wolf—and lone hobgoblins are \textbf{more dangerous}, because they have nothing left to lose.

\textbf{Resolution}: Hobgoblins can be reasoned with, but only from a position of strength:
\begin{itemize}
    \item \textbf{Earn Respect}: Succeed on a Combat or Command test (DV 4) in front of the hobgoblin. On success, negotiate terms at Controlled Position.
    \item \textbf{Pay in Service}: Offer to perform a task or service in exchange for passage or aid. The GM sets a Progress Timer (4—6 segments) representing the task.
    \item \textbf{Oath of Blood}: Spend 1 Boon to bind a temporary oath. The hobgoblin will not attack you or your allies for one scene.
\end{itemize}

\section{Bugbears — Threshold's Stalker}

\textbf{Description}: Large, silent, and patient. Bugbears are the hunters of the goblinoid kind. They move with a quiet that seems impossible for their size, their fur matted and dark, their eyes reflecting the dimmest light. They do not speak often. They do not need to.

Bugbears are the Threshold's \textbf{hunters}. They appear where something has been lost—a name, a memory, a child. They do not find. They \textbf{track}. And they do not stop until the balance is paid.

\subsection{Bugbear NPC Template}

\textbf{Tier}: III | \textbf{Cap}: 4 | \textbf{Corporeal}, Silent, Patient

\textbf{Dice Pools}:
\begin{itemize}
    \item \textbf{Melee}: 8d (Body 4 + Melee 4) — \textit{Harm 5 (claws or heavy weapon)}
    \item \textbf{Stealth}: 7d (Wits 4 + Stealth 3)
    \item \textbf{Athletics}: 7d (Body 4 + Athletics 3)
    \item \textbf{Endurance}: 7d (Body 4 + Endurance 3)
\end{itemize}

\textbf{Armor}: Light (thick fur) — \textit{Harm 1 → Fatigue 1}

\textbf{Traits}:
\begin{itemize}
    \item \textbf{Threshold's Patience} — Once per scene, the bugbear may treat a failed Stealth roll as a Clean Success. The GM gains 1 SB.
    \item \textbf{Surprise Strike} — If the bugbear attacks from Hidden, it deals \textbf{+2 Harm} on the first hit.
    \item \textbf{Unnatural Stillness} — The bugbear can hold a position, unmoving, for hours. It cannot be detected by any mundane sense while holding still.
\end{itemize}

\textbf{Etiquette}: Bugbears do not bargain. They do not negotiate. They hunt. The only question is whether they consider you prey or \textbf{interesting}. A bugbear that is curious will watch. A bugbear that is hungry will strike. A bugbear that is bored will leave. The key is to be interesting enough to watch, but not interesting enough to hunt.

\textbf{Complications}: Bugbears are loners by nature, but they will occasionally form packs. A pack of bugbears is a death sentence in the wild. They do not retreat. They do not negotiate. They simply \textbf{pursue}.

\textbf{Resolution}: A bugbear can be distracted by a more interesting target, or by something that requires its patience to observe:
\begin{itemize}
    \item \textbf{Offer a Secret}: Spend 1 Boon to offer a secret, memory, or name. The bugbear's attention shifts, granting the party a window to escape or act. The GM gains 1 SB as the bugbear files the information.
    \item \textbf{Be Interesting}: Succeed on a Performance, Lore, or Deception test (DV 4) to capture the bugbear's curiosity. On a Clean Success, the bugbear watches but does not attack for one scene. On a Partial, it watches—but with hunger.
    \item \textbf{Provide a Greater Hunt}: Point the bugbear toward a more worthy target (a rival, a monster, a dangerous individual). Succeed on a Sway or Deception test (DV 4). On success, the bugbear leaves you alone to pursue the new quarry.
\end{itemize}

\section{Wargs — Threshold's Howl}

\textbf{Description}: Lean, wolf-like predators with matted fur and eyes that gleam like wet stone. Wargs move in loose packs along threshold edges, drawn to places where balance has frayed. Their howls carry no malice—only the sharp insistence of unresolved tension. They do not hunt for sport, but to remind the world of what it owes itself.

Wargs are the Threshold's \textbf{reminders}. They appear when something has been overlooked—a debt unpaid, a promise broken, a name forgotten. They do not attack. They \textbf{watch}. And their watching is a pressure that builds until the balance is restored.

\subsection{Warg NPC Template}

\textbf{Tier}: II—IV | \textbf{Cap}: 3 | \textbf{Corporeal}, Persistent, Attuned

\textbf{Dice Pools}:
\begin{itemize}
    \item \textbf{Melee}: 6d (Body 3 + Brawl 3) — \textit{Harm 3 (bite or claw)}
    \item \textbf{Athletics}: 6d (Body 3 + Athletics 3)
    \item \textbf{Notice}: 6d (Wits 3 + Notice 3)
    \item \textbf{Survival}: 6d (Wits 3 + Survival 3)
\end{itemize}

\textbf{Armor}: Light (thick hide) — \textit{Harm 1 → Fatigue 1}

\textbf{Traits}:
\begin{itemize}
    \item \textbf{Threshold's Echo} — Once per scene, the warg may reroll a failed Notice or Survival roll when tracking imbalance. The GM gains 1 SB.
    \item \textbf{Pack's Patience} — As a Reaction when the pack would take Harm, it may ignore 1 Fatigue. The GM gains 1 SB.
    \item \textbf{Howl of Unpaid Balance} — Once per scene, a warg pack's howl forces all who hear it to test Spirit + Resolve (DV 3) or gain 1 Fatigue as the weight of neglected duty presses upon them. The GM gains 1 SB for each character who fails.
\end{itemize}

\textbf{Etiquette}: Wargs respect quiet endurance and unflinching presence. A character who stands firm without aggression gains \textbf{+1 Position} when encountered. They ignore those who flee or overreact—only the steady draw their attention.

\textbf{Complications}: Warg packs dissolve if their perceived imbalance is corrected—but reform swiftly if new tension arises. They bear no grudges, only acute awareness of what remains unsettled.

\textbf{Resolution}: Most wargs withdraw when offered patient observation or a shared moment of silence:
\begin{itemize}
    \item \textbf{Match Stillness}: Succeed on a Spirit + Endurance test (DV 3). On a Clean Success, the wargs withdraw or linger as silent guardians. On a Partial, they withdraw but the GM gains 1 SB.
    \item \textbf{Share Silence}: Spend 1 Boon to offer a moment of stillness. The wargs observe but do not attack for one scene.
    \item \textbf{Acknowledge Imbalance}: Succeed on a Lore or Insight test (DV 4) to name the imbalance the wargs track. On success, they withdraw. On failure, the GM gains 1 SB and the wargs press closer.
\end{itemize}

\section{Goblinoid Bargains — Threshold's Interest}

All goblinoids understand one thing: \textbf{balance}. A goblin who is given a coin will remember. A hobgoblin who is spared will owe. A bugbear who is not hunted will wonder. The key to dealing with goblinoids is to \textbf{never owe them more than you can pay}.

\subsection{Payment Options}

\begin{table}[h]
\centering
\begin{tabular}{|l|l|l|}
\hline
\textbf{Payment Type} & \textbf{Effect} & \textbf{Cost} \\
\hline
\textbf{Food} & Goblins accept food for service. The food must be fresh, warm, and shared. & Spend 1 Supply or offer a meal \\
\textbf{Shiny Objects} & Goblins are attracted to anything that catches the light. A coin, a broken mirror, a shard of glass—all are currency. & Spend 1 Boon or offer a valuable trinket \\
\textbf{Names} & Hobgoblins value names. A name given freely is a name that can be called. A name that is called is a name that can be sold. & Offer a name (grants GM 1 String: The Name Known) \\
\textbf{Silence} & Bugbears value silence. A secret kept, a truth withheld, a story never told—these are the prices they accept. & Offer a secret (grants GM 1 SB as the secret enters circulation) \\
\hline
\end{tabular}
\caption{Goblinoid Payment Options}
\end{table}

\subsection{Bargain Resolution}

\begin{enumerate}
    \item \textbf{Offer Payment}: Declare what you offer (food, shiny object, name, silence).
    \item \textbf{Set Position}: The GM sets Position based on the goblinoid's mood and the value of the offer (typically Controlled to Desperate).
    \item \textbf{Roll}: Sway (Presence + Sway) for negotiation, or appropriate skill (Insight, Deception, Command).
    \item \textbf{Outcome}:
    \begin{itemize}
        \item \textbf{Clean Success}: The bargain is accepted. The goblinoid provides service or leaves peacefully.
        \item \textbf{Partial}: The bargain is accepted, but the GM gains 1 SB (a hidden condition or future complication).
        \item \textbf{Miss}: The goblinoid refuses. It may become hostile or demand a higher price. The player gains 2 Boons.
    \end{itemize}
\end{enumerate}

\subsection{When Payment Fails}

If you cannot pay or refuse to pay, the goblinoids will collect. The Threshold's interest, after all, is always compounding. The GM may:

\begin{itemize}
    \item Spend 2 SB to call in a debt during a crucial moment.
    \item Have a goblinoid appear at an inopportune moment to collect—at Desperate Position.
\end{itemize}

\section{Goblinoid Encounters — Framing and Resolution}

\subsection{Encounter Setup}

\begin{enumerate}
    \item \textbf{Determine the Goblinoid's Goal}: What do they want? (food, passage, payment, information, a name).
    \item \textbf{Set Initial Position}: Based on the party's recent actions and the goblinoid's mood.
    \begin{itemize}
        \item \textbf{Dominant}: The party has leverage (intimidated a goblin, earned respect).
        \item \textbf{Controlled}: Neutral encounter (the goblinoids are watching, waiting).
        \item \textbf{Desperate}: The party is cornered, outnumbered, or has failed a bargain.
    \end{itemize}
\end{enumerate}

\subsection{Encounter Resolution}

\begin{itemize}
    \item \textbf{Social Resolution}: Use the Bargain mechanics above. Sway, Command, or Deception.
    \item \textbf{Combat Resolution}: Use standard Fate's Edge combat. Goblins (Cap 1), Hobgoblins (Cap 3), Bugbears (Cap 4).
    \item \textbf{Clever Resolution}: A well-placed trick or clever idea may bypass the encounter entirely. The GM may set a DV 3—4 test and allow the party to spend Boons to improve Position.
\end{itemize}

\subsection{GM Quick Cues}

\begin{itemize}
    \item Goblins are \textbf{opportunistic}: they will bargain, steal, or flee depending on the situation. They are \textbf{fodder}—do not overcomplicate them.
    \item Hobgoblins are \textbf{tactical}: they will not attack unless they believe they can win or until their patience is exhausted.
    \item Bugbears are \textbf{patient}: they will watch and learn before acting. They are a threat that grows with time.
    \item Wargs are \textbf{attuned}: they respond to imbalance, not aggression. A fight with wargs is a failure of awareness, not a test of strength.
\end{itemize}

\section{Quick Reference — Goblinoids}

\begin{table}[h]
\centering
\begin{tabular}{|l|l|l|l|l|l|l|}
\hline
\textbf{Creature} & \textbf{Tier} & \textbf{Cap} & \textbf{Key Traits} & \textbf{Harm} & \textbf{Armor} & \textbf{Resolution} \\
\hline
\textbf{Goblin} & I—II & 1 & Threshold's Cunning, Nimble Escape & 2 & None & Bribe, trick, or outsmart \\
\textbf{Tunnel-Goblin} & I & 1 & Echolocation, Skittering Charge & 2 & None & Light, or outrun \\
\textbf{Ember-Goblin} & I & 1 & Smoldering Touch, Ember Burst & 2 & None & Water, or kill at range \\
\textbf{Sea-Goblin} & I & 1 & Amphibious, Drowning Grasp & 2 & None & Stay away from water \\
\textbf{Star-Goblin} & I & 1 & Star-Struck, Starlight Shimmer & 2 & None & Kill or ignore \\
\textbf{Hobgoblin} & II—III & 3 & Threshold's Discipline, Tactical Awareness & 4 & Medium (1→1, 3→2) & Earn respect, pay in service \\
\textbf{Bugbear} & III & 4 & Threshold's Patience, Surprise Strike & 5 & Light (1→1) & Offer a secret, be interesting \\
\textbf{Warg} & II—IV & 3 & Threshold's Echo, Pack's Patience & 3 & Light (1→1) & Match stillness, acknowledge imbalance \\
\hline
\end{tabular}
\caption{Goblinoid Quick Reference}
\end{table}

\section{Goblinoid Spawn — Threshold's Progeny (Optional)}

When the Threshold is particularly thin or agitated, goblinoids may spawn spontaneously from the friction of broken promises, forgotten debts, or unresolved tensions. These spawn are not born—they are \textbf{manifested}.

\subsection{Spawn Mechanics}

When a goblinoid spawns, the GM may:

\begin{itemize}
    \item \textbf{Spend 1 SB} to manifest a single goblin from a broken promise or forgotten debt.
    \item \textbf{Spend 2 SB} to manifest a hobgoblin from a broken oath or an unbalanced transaction.
    \item \textbf{Spend 3 SB} to manifest a bugbear from a forgotten grudge or a long-suppressed truth.
    \item \textbf{Spend 4 SB} to manifest a warg pack from a neglected duty or a frayed threshold.
\end{itemize}

Each manifestation:
\begin{itemize}
    \item Occurs at the most inconvenient moment (GM's choice).
    \item The party may resolve the spawn by addressing the underlying imbalance.
\end{itemize}

\subsection{Spawn Examples}

\begin{table}[h]
\centering
\begin{tabular}{|l|l|l|}
\hline
\textbf{Trigger} & \textbf{Manifestation} & \textbf{Resolution} \\
\hline
A broken promise & Goblin (Cap 1) arrives to collect the unpaid debt & Fulfill the promise or pay in service \\
A forgotten oath & Hobgoblin (Cap 3) appears at the scene of the betrayal & Restore the oath or offer a name \\
A suppressed truth & Bugbear (Cap 4) begins tracking those who knew the secret & Reveal the truth or offer a greater secret \\
A frayed threshold & Wargs (Cap 3, pack of 2—3) begin patrolling the area & Acknowledge the imbalance and restore the threshold \\
\hline
\end{tabular}
\caption{Goblinoid Spawn Examples}
\end{table}

\section{GM Notes: Running Goblinoids in Fate's Edge}

\subsection{Narrative Role}

Goblinoids serve as the Threshold's \textbf{interest collectors}. They appear when:
\begin{itemize}
    \item A promise has been broken.
    \item A debt has gone unpaid.
    \item A balance has been disturbed.
\end{itemize}

They are not random encounters—they are \textbf{consequences}. When the party fails a social roll, ignores a bond, or leaves a thread unresolved, the GM may use goblinoids to introduce pressure and reinforce the theme of balance.

\subsection{Using the Goblinoid Toolkit}

\begin{enumerate}
    \item \textbf{Manifestation}: When the fiction demands it, a goblinoid appears to collect. Use the appropriate template.
    \item \textbf{Bargain}: Offer the party a chance to resolve the situation through bargaining, cleverness, or payment.
    \item \textbf{Escalation}: If the party refuses, the goblinoid attacks, steals, or calls in a debt. The GM gains SB to escalate further.
\end{enumerate}

\subsection{Integration with Campaign Themes}

Goblinoids are ideal for:
\begin{itemize}
    \item \textbf{Bond-Driven Play}: A goblinoid may appear to collect a debt tied to a character's Bond.
    \item \textbf{Complication Spotlight}: A character's Complication (Debt to a Patron, Hunted, Oath-Bound) may manifest as a goblinoid encounter.
    \item \textbf{Travel Framework}: Wargs are excellent encounters during travel legs, especially when the party has ignored threshold markers or rushed past significant sites.
    \item \textbf{Magic System}: Runekeepers and Invokers who push their Obligation may attract goblinoid attention—the Threshold's interest is patient, but it does not forget.
\end{itemize}

\subsection{Final Thought}

The goblinoids are not evil. They are not monsters. They are the price of forgetting what you owe—the Threshold's reminder that every action has a cost, and every cost must be paid. They will collect. The only question is whether you pay willingly, or whether they take what they are owed.

\lerris{I will not write more. Some debts are not meant to be recorded. But I have seen the goblinoids' eyes when they think no one is watching. They are not afraid. They are \textbf{waiting}. For what, I do not know. I do not think I want to find out.}